\begin{table}[t]
\centering
\caption{What the HEALPix$\to$UTM read does beyond preserving a constant}
\label{tab:operator_tests}
\footnotesize
\begin{tabular}{@{}rccccr@{}}
\toprule
Wavelength & \multicolumn{4}{c}{Amplitude response} & Max shift \\
\cmidrule(lr){2-5}
(m) & $0^\circ$ & $45^\circ$ & $90^\circ$ & $135^\circ$ & (m) \\
\midrule
12 & 0.568 & 0.810 & 0.514 & 0.353 & 0.00 \\
21 & 0.837 & 0.933 & 0.806 & 0.732 & 0.00 \\
37 & 0.943 & 0.977 & 0.930 & 0.902 & 0.00 \\
63 & 0.980 & 0.992 & 0.976 & 0.966 & 0.00 \\
107 & 0.993 & 0.997 & 0.992 & 0.988 & 0.00 \\
184 & 0.998 & 0.999 & 0.997 & 0.996 & 0.00 \\
316 & 0.999 & 1.000 & 0.999 & 0.999 & 0.00 \\
541 & 1.000 & 1.000 & 1.000 & 1.000 & 0.00 \\
928 & 1.000 & 1.000 & 1.000 & 1.000 & 0.00 \\
1592 & 1.000 & 1.000 & 1.000 & 1.000 & 0.00 \\
\bottomrule
\end{tabular}
\vspace{3pt}
\begin{minipage}{\columnwidth}
\footnotesize
Each row is a single-frequency wave placed analytically on the level-20 cell centres, read on the native UTM pixels, and regressed against the same wave evaluated at those pixels; the columns are the recovered amplitude response by wave direction in the projected frame, and the largest phase shift converted to metres. Nothing is resampled: the ground truth is the same analytic function, so the residual is the operator and the unavoidable discretisation of reading a cell centre. On a band-limited field with the imagery's own fitted spectrum the operator contributes 1.4e-01 of the field's standard deviation; one cell lights 1.0 native pixels; a unit gradient is reproduced to 0.32 cells.
\end{minipage}
\end{table}
